\chapter{Countermodels, Boundary Cases, and Worked Constructions}
\label{sec:p14-countermodels-and-boundary-cases}
\index{countermodel!reference}\index{counterexample!hypothesis dependence}

A countermodel to an inference satisfies every premise and fails its
conclusion under the specified semantics. A boundary construction can instead
show that a theorem's hypotheses cannot simply be removed, that an object
does not exist in a given class, or that observations leave a target
undetermined. These uses require different witnesses. The constructions below
provide explicit objects and calculations for consultation, with links to
their governing definitions and proof details.

The examples concern their named languages, structures and procedures. A
finite valuation refutes one proposed classical inference; a family of
finite structures indexed by every rank is needed for a first-order
inexpressibility argument. An impossibility theorem quantifies over all
eligible procedures or models, rather than reporting that one attempted
procedure failed. Sources and proof standing remain those of the linked
entries \cite{Enderton2001,Libkin2004,Kleene1952}.

\section{Small logical models and separating valuations}
\label{sec:p14-counter-small-models}

\subsection{A two-value valuation}
\label{sec:p14-counter-boolean}

Let \(p,q\) be distinct propositional atoms and let a classical valuation
assign \(v(p)=\mathbf F\), \(v(q)=\mathbf T\). Using the truth
clauses of \cref{sec:prop-syntax-semantics-02} gives:
\[
\begin{array}{c|cccccc}
 \text{formula}&p&q&p\to q&\neg p&\neg q&p\land q\\\hline
 \text{value}&\mathbf F&\mathbf T&\mathbf T&\mathbf T&\mathbf F&\mathbf F
\end{array}
\]
Consequently \(p\to q,q\) does not entail \(p\), and
\(p\to q,\neg p\) does not entail \(\neg q\). The very same assignment
satisfies each pair of premises and falsifies its proposed conclusion.
These are the converse and inverse failures in
\cref{def:p18-rd-consequent}, verified with the related reconstructions in
\cref{thm:p18-rd-independent-cases}. They leave valid contraposition
\(p\to q,\neg q\entails\neg p\) intact.

\subsection{A nonempty universe with an empty predicate}
\label{sec:p14-counter-empty-predicate}

Take the first-order domain \(D=\{s\}\), and interpret unary predicates
oak, tree and plant by the empty set. Every oak is a tree and every tree
is a plant are true, as is every oak is an oak. Some oak is a plant and
some oak is an oak are false. The universe has an element, but none
satisfies oak. This is the complete model of
\cref{def:p02-syllogistic-empty-term}; its categorical formulas use the
semantics of \cref{sec:p02-syllogistic-inference}
\cite[sec.~3]{PrattHartmannMoss2008SyllogisticV1}.

A separate one-element model with the first predicate empty and the second
equal to \(D\) satisfies every first-predicate object is a second-predicate
object, while falsifying the converse. The empty-subject and converse
failures both preserve nonempty-domain semantics. An empty universe is a
different boundary: \cref{def:p02-syllogistic-nonempty-universe} explains
why the literal-complement rule N uses the nonemptiness assumption of
its exact \(\mathsf S^\dagger\) presentation.

\subsection{A persistent intuitionistic model}
\label{sec:p14-counter-ipl}

Use the two-element chain \(w_0<w_1\), with the reflexive order pairs
included, and the persistent atomic interpretation \(V(p)=\{w_1\}\).
No world forces falsity. The exact implication and negation clauses are
\cref{def:p08-ipl-forcing}. At \(w_0\), \(p\) is not forced. Neither
is \(\neg p\), because the accessible \(w_1\) forces \(p\).
Thus \(p\lor\neg p\) is not forced at \(w_0\): intuitionistic
disjunction requires a disjunct there, not eventual satisfaction of one
at a later point. At \(w_1\), \(p\) is forced and \(\neg p\) is not.

Neither world forces \(\neg p\); hence \(w_0\) forces \(\neg\neg p\)
while it does not force \(p\). This same model separates double-negation
elimination from the intuitionistic semantics. The model and its excluded-
middle evaluation are \cref{fig:p08-ipl-excluded-middle}, under
\cref{def:p08-ipl-kripke-structure}. The separately displayed Peirce
countermodel is \cref{fig:p08-ipl-peirce}
\cite{Kripke1965,TroelstraVanDalen1988}. A classical two-value assignment
cannot replace this persistence model when testing intuitionistic validity.

\subsection{The LP designation boundary}
\label{sec:p14-counter-lp}

In the exact LP matrix, assign \(A\) the value both and \(B\) the value
false. Both \(A\) and \(\neg A\) are designated, whereas \(B\) is not.
This refutes designation-preserving consequence
\(\{A,\neg A\}\entails_{\mathsf{LP}}B\), as in
\cref{def:p08-lp-nonexplosion}. LP uses the strong-Kleene operations and
designates the values true and both; the different designated set of K3 gives
a different consequence relation even with the same operations
\cite{Priest1979}. This valuation establishes nonexplosion. It does not
assert that every contradictory theory is nontrivial or that a contradiction
describes an actual situation.

\subsection{Finite modal frames}
\label{sec:p14-counter-modal}

For each row of the following table, \(W\) is the listed nonempty world
set, \(R\) contains exactly the listed pairs, and \(V(p)\) is the listed
set of worlds where the atom is true. The witness is evaluated at the
named world. The box clause quantifies over its \(R\)-successors and
\(\Diamond p\) abbreviates \(\neg\Box\neg p\).

\begin{longtable}{@{}p{.19\textwidth}p{.38\textwidth}p{.33\textwidth}@{}}
\caption{Finite modal frames}\\
\toprule
Target and world & Complete frame and valuation & Evaluated failure \\
\midrule
\endfirsthead
\caption[]{Finite modal frames (continued)}\\
\toprule
Target and world & Complete frame and valuation & Evaluated failure \\
\midrule
\endhead
\midrule\multicolumn{3}{r}{\emph{continued on next page}}\\
\endfoot
\endlastfoot
T in \(KD\), at \(w\) &
\(W=\{w,v\}\), \(R=\{(w,v),(v,v)\}\), \(V(p)=\{v\}\). &
The frame is serial; \(\Box p\) holds at \(w\), but \(p\) does not. \\
D in \(K4\), at \(w\) &
\(W=\{w\}\), \(R=\varnothing\), any \(V(p)\). &
The frame is transitive; \(\Box p\) is vacuously true and
\(\Diamond p\) is false. \\
4 in \(KTB\), at \(a\) &
\(W=\{a,b,c\}\); all loops and \((a,b),(b,a),(b,c),(c,b)\);
\(V(p)=\{a,b\}\). &
The frame is reflexive and symmetric. Both successors of \(a\) satisfy
\(p\), but successor \(b\) sees \(c\), so \(\Box\Box p\) fails at \(a\). \\
B in \(K4\), at \(a\) &
\(W=\{a,b\}\), \(R=\{(a,a),(a,b),(b,b)\}\), \(V(p)=\{a\}\). &
The frame is transitive. At \(a\), \(p\) holds but successor \(b\)
has no \(p\)-successor, falsifying \(\Box\Diamond p\). \\
T in symmetry-only \(KB\), at \(a\) &
\(W=\{a,b\}\), \(R=\{(a,b),(b,a)\}\), \(V(p)=\{b\}\). &
Symmetry holds without reflexivity; \(\Box p\) is true and \(p\) false
at \(a\). \\
\bottomrule
\end{longtable}

These are the complete witnesses in
\cref{thm:p08-mce-separations,annex:p8.mce.10}. Their relevant Hilbert
sheets are sound on the stated frame classes, so a failed axiom instance
also establishes the indicated nonderivability. In particular, the
symmetry-only \(KB\) row is not a countermodel on the reflexive symmetric
frames of \(KTB\). The source and notation conventions are
\cref{def:p08-mce-frame-classes}\cite{Chellas1980}.

\section{Counterexamples to proposed inference rules}
\label{sec:p14-counter-inference-rules}

\subsection{Converse, inverse and distribution of a negative answer}

The valuation in \cref{sec:p14-counter-boolean} distinguishes an implication
from its converse or inverse. Adding a justified converse premise
\(q\to p\) would repair the converse deduction; treating \(p\) as a
possible explanation of \(q\) asks a different, abductive question.
The observation alone does not establish that explanatory hypothesis.

For the proposed inference
\(\neg(p\land q)\entails\neg p\land\neg q\), instead put
\(p=\mathbf T,q=\mathbf F\). The premise is true but the conclusion
is false. The valid classical consequence is the disjunction
\(\neg p\lor\neg q\). This is the Boolean construction behind
\cref{def:p18-rd-compound-question}; a negative answer to whether both
individuals have a property need not deny it of both. The proof of these
modern diagnostic reconstructions is
\cref{thm:p18-rd-independent-cases}. The historical classification retains
the selected Pickard-Cambridge translation's separate role
\cite{AristotleSophisticalRefutationsPickardCambridgeMIT}.

\subsection{Dropping an eigenvariable condition}
\label{sec:p14-counter-eigenvariable}

The initial sequent \(P(y)\Rightarrow P(y)\) does not license
\(P(y)\Rightarrow\forall xP(x)\) by an unrestricted \(\forall R\).
Take domain \(\{0,1\}\), predicate interpretation \(P=\{0\}\), and
assignment \(y\mapsto0\). The first sequent is true, and the antecedent
of the proposed conclusion is true. Its universal succedent is false at
the domain element 1. Thus allowing that inference without the freshness
condition would be unsound.

This is \cref{def:p07-classical-eigenvariable-case}. The defined-negation
classical calculus is fixed on \cref{sheet:lk-defneg}; it is distinct from
the primitive-negation \cref{sheet:lk} presentation. Renaming an
eigenvariable means coherently renaming the premise derivation while retaining
its side conditions. Merely replacing its printed name does not allow the
variable to remain free in a forbidden context.

\subsection{Necessitating a local premise}
\label{sec:p14-counter-local-necessitation}

Take \(W=\{a,b\}\), the full relation \(R=W^2\), and
\(V(p)=\{a\}\). The frame is serial, reflexive, symmetric and transitive.
At \(a\), premise \(p\) holds, but \(\Box p\) fails because \(b\)
is accessible and does not satisfy \(p\). Consequently
\(\{p\}\not\derives_L\Box p\) for each locally sound
\(L\in\{KD,KTB,K4\}\).

If \(p\) is instead assumed true at every world, every successor of every
world satisfies it and \(\Box p\) is true globally. The change is in
the consequence relation's premise condition. The complete comparison and
proof are \cref{thm:p08-mce-separations,annex:p8.mce.10}. The
theorem-only necessitation convention of \cref{def:p08-mce-language}
preserves local soundness; the global reading does not justify inserting
premise-dependent necessitation into that local calculus.

\subsection{Existence, maximality and empty families}
\label{sec:p14-counter-extensions}

In the one-argument attack framework \((\{a\},\{(a,a)\})\), the
set \(\{a\}\) is not conflict-free. The empty set is its grounded,
complete and preferred extension. It has no stable extension: the empty
set attacks no outside argument, and the singleton cannot be conflict-free.
The characteristic operator nevertheless satisfies
\(\mathsf F(\{a\})=\{a\}\), since the member attacks its own sole
attacker. Thus an arbitrary fixed point of the defense operator need not
be a complete extension.

For the empty framework \((\varnothing,\varnothing)\), by contrast,
\(\varnothing\) is a stable extension, and the stable-extension family
is \(\{\varnothing\}\). A family with no members and a family containing
the empty extension are different. The guarded skeptical-acceptance query
of \cref{def:p18-acceptance-queries} requires that an extension exist;
it reports no extensions for the first framework's stable semantics.
The complete lists and proofs are
\cref{def:p18-example-frameworks,thm:p18-example-extensions}
\cite{Dung1995}.

The unequal-size framework in the same entry, with attacks
\(a\leftrightarrow b\) and \(a\leftrightarrow c\) and no others,
has preferred extensions \(\{a\}\) and \(\{b,c\}\). Both are maximal
under inclusion, despite their different sizes. Replacing maximal by
maximum-cardinality would change the semantics, not clarify its definition.

\section{Failures of compactness, completeness, definability, and decidability}
\label{sec:p14-counter-metatheory}

\subsection{Finite satisfiability without a finite model of the whole theory}
\label{sec:p14-counter-finite-compactness}

Use a first-order signature with a strict-order relation \(<\) and equality.
Let the theory contain the strict-linear-order axioms and, for each integer
\(m\geq1\), the sentence
\[
 \exists x_1\cdots\exists x_m\,
                  \bigwedge_{1\leq i<j\leq m}x_i\ne x_j.
\]
For \(m=1\) the empty conjunction is true. Every finite portion of these
sentences has a largest required number of distinct objects; a finite
linear order at least that large satisfies the portion. No finite domain
satisfies all the sentences. This is
\cref{def:p03-finite-compactness-failure}. It shows failure of compactness
under finite-model semantics. Ordinary first-order compactness permits an
infinite order and therefore is not contradicted
\cite{Libkin2004,EbbinghausFlum1995}.

\subsection{A full second-order theory with only finite-subset models}
\label{sec:p14-counter-second-order-compactness}

Use the full second-order \(\mathsf{PA}_2\) of \cref{def:p02-pa2},
extend its signature by a constant \(c\), and form
\[
 T=\{\mathsf{PA}_2\}\cup\{c\ne S^n0:n\in\N\}.
\]
For a finite subset, use the standard natural-number structure and choose
\(c\) different from the finitely many excluded numerals. The full
second-order induction sentence is satisfied because its relation quantifier
ranges over every subset of the domain. A full model of the whole theory
would be isomorphic to the standard natural numbers by
\cref{thm:p02-pa2-categoricity}, yet would contain \(c\) different from
every numeral. That is impossible. The complete proof accompanies
\cref{thm:p02-sol-full-noncompact} in
Section~\ref{sec:second-order-logic-04}.

This witness depends on full relation ranges. It does not transfer to Henkin
semantics, whose permitted relation objects and comprehension conditions are
part of a different model class. Nor does noncompactness alone prove that
every notion of effective theoremhood fails. The distinct entry
\cref{thm:p02-sol-full-no-effective-completeness} specifies computably
checkable finite derivations and reduces true first-order arithmetic to full
validity. Its inline argument in Section~\ref{sec:second-order-logic-04}
is labelled a proof sketch, with the
effective coding and non-enumerability details left to the cited sources
\cite{BarwiseFeferman1985,Shapiro1991}. The compactness witness proves
noncompactness; it does not provide those effective-coding and arithmetic
non-enumerability arguments.

\subsection{A structure pair for every first-order rank}
\label{sec:p14-counter-definability}

For each integer \(k\geq1\), take the finite strict linear orders of sizes
\(2^k\) and \(2^k+1\). One has even cardinality and the other odd.
The strategy in \cref{def:p03-ef-parity} matches selected points in order,
maintaining that corresponding gaps, with \(r\) rounds remaining, either
have equal size below \(2^r-1\) or both have at least that size.
The gap size counts unselected elements between successive selected points,
including the two end gaps. A repeated move receives its existing mate.

In equal small gaps, copy the exact offset. In two large gaps, copy a side
smaller than \(2^{r-1}-1\) if there is one; otherwise choose an answer
with at least \(2^{r-1}-1\) elements on either side. Both possibilities
preserve the invariant for the next round. At zero rounds the pairing
preserves order and equality. The rank-zero case is vacuous. Thus there is
an even/odd indistinguishable pair for every finite rank, meeting the
quantifier in \cref{thm:p03-ef-inexpressibility}. No first-order sentence
defines even cardinality on finite linear orders.

A single pair agreeing through one chosen rank would only exclude formulas
up to that rank. It would not prove this unrestricted inexpressibility
statement. The game theorem and criterion at
\cref{thm:p03-ef,thm:p03-ef-inexpressibility} retain their cited standing
\cite{Ehrenfeucht1961,Libkin2004}.

For finite graph connectivity, the object pair is instead one cycle
\(C_{2m}\) versus the disjoint union \(C_m\sqcup C_m\), with \(m\)
large enough for the selected rank. The first is connected and the second
is not. \Cref{def:p03-ef-connectivity} records the bounded-neighborhood
strategy and cites its full invariant and response lemma to Libkin. The
parity proof is not a substitute for that graph argument; the sufficient
cycle size and available responses remain conditions of the cited construction.

\subsection{A diagonal counterexample to every proposed domain decider}
\label{sec:p14-counter-decidability}

Fix the program-numbering and application conventions of
Section~\ref{sec:universal-reductions-01} and
the self-application domain set
\(K=\{e:\varphi_e(e)\text{ is defined}\}\).
Defined means acceptance with a valid numerical output under that numbering,
not any physical halt. Suppose a total computable characteristic function
\(h\) decided \(K\), using value one for membership. Define the partial
procedure \(d(e)\) to return zero when \(h(e)=0\) and otherwise diverge.
The enumeration theorem supplies an index \(k\) for \(d\).
If \(h(k)=0\), then \(\varphi_k(k)=0\), contradicting nonmembership.
If \(h(k)=1\), then \(\varphi_k(k)\) is undefined, contradicting membership.
Both answers of a purported decider fail on its associated diagonal index.

The complete construction and proof accompany
\cref{thm:halting-undecidable}, using
\cref{thm:enumeration}\cite{Kleene1952}. Recognition is different:
simulating until a valid output appears recognizes membership without
necessarily terminating on nonmembers. The theorem supplies an impossibility
argument for a total decider; an individual unsuccessful program search does
not. Likewise the finite-validity result \cref{thm:trakhtenbrot} retains
its finite relational signature with a binary relation. That theorem does
not classify every restricted fragment or the problem of evaluating one
explicit finite structure \cite{Trakhtenbrot1950}.

\section{Statistical, causal, decision, and aggregation counterexamples}
\label{sec:p14-counter-statistical-causal-decision}

\subsection{Pairwise independence and selection}
\label{sec:p14-counter-independence}

Give each outcome in \(\Omega=\{(0,0),(0,1),(1,0),(1,1)\}\) mass
\(1/4\). Let events \(A,B\) mean that the first and second coordinates,
respectively, are one, and let \(C\) mean that the coordinates agree.
Write \(P\) for the event probability induced by these four weights.
All three probabilities are \(1/2\), and every pairwise intersection
has probability \(1/4\). Nevertheless
\[
 P(A\cap B\cap C)=1/4\ne1/8=P(A)P(B)P(C).
\]
Thus pairwise independence does not imply mutual independence. Conditional
on \(C\), the two retained outcomes are equally weighted, so
\[
 P(A\mid C)=P(B\mid C)=1/2,\qquad P(A\cap B\mid C)=1/2.
\]
Selection has made the initially independent coordinates dependent. This
is \cref{def:p15-pairwise-not-mutual}, under
\cref{def:p15-finite-independence,def:p15-finite-conditioning}. The stronger
fifteen-variable counterexample in \cref{thm:p15-conc-dependence-boundaries}
also exhibits failure of an independent-sum concentration bound when mutual
independence is replaced by pairwise independence.

\subsection{A likelihood with a different posterior}
\label{sec:p14-counter-base-rate}

The stipulated defect/alert model of \cref{def:p15-alert-bayes-case} has
the following joint weights, divided by \(10000\):
\[
\begin{array}{c|rr}
 &\text{alert}&\text{no alert}\\\hline
 \text{defect}&90&10\\
 \text{no defect}&495&9405
\end{array}
\]
It gives \(P(\text{alert}\mid\text{defect})=9/10\),
\(P(\text{alert}\mid\text{no defect})=1/20\), and defect prior
\(1/100\). Among the alert weight \(585\), only \(90\) is defective:
\[
 P(\text{defect}\mid\text{alert})=90/585=2/13.
\]
The likelihood \(9/10\) is not this posterior. Holding the two likelihoods
fixed but changing the prior to \(1/10\) gives posterior \(2/3\).
The same source entry also calculates genuinely conditionally independent
repeated alerts. Duplicating one report requires the separate dependence
calculation in \cref{thm:p15-bayes-flag-calculation}; multiplying the same
likelihood ratio twice does not represent a copied observation.

\subsection{Posterior content without uniform sampling coverage}
\label{sec:p14-counter-coverage}

For one Bernoulli observation with a uniform prior on its parameter, the
equal-tail posterior intervals of content \(0.95\) are
\[
 C(1)=[\sqrt{0.025},\sqrt{0.975}],\qquad
 C(0)=[1-\sqrt{0.975},1-\sqrt{0.025}].
\]
Both exclude the fixed parameter \(0.99\), since their upper endpoints
are approximately \(0.9874\) and \(0.8419\). The coverage probability
at that parameter is therefore zero, even though either observation gives
the stated posterior content. \Cref{thm:p15-bayes-coverage} supplies the
complete proof and the separate prior-average calibration identity.
The distinction is between a conditional parameter law and repeated data
at a fixed parameter, as defined in \cref{def:p15-bayes-credible-set}.
It is not a contradiction between those definitions.

\subsection{A perfect sample fit selected from many candidates}
\label{sec:p14-counter-selection}

Fix integers \(K,n\geq1\). Let every record have \(K\) independent fair
binary features and a label always equal to zero. The fixed candidate
predictor \(h_j\), \(1\leq j\leq K\), reports feature \(j\).
Each candidate has population error \(1/2\). On \(n\) independent
records, a particular predictor has empirical error zero with probability
\(2^{-n}\). Independence across the feature columns gives probability
\[
 1-(1-2^{-n})^K
\]
that at least one does so. This can be large although no candidate has
small population error. The exact construction and proof are
\cref{thm:p15-conc-selection-example,annex:p15.conc-selection-example}.
Using one fixed-candidate error bound after searching the entire collection
changes its target. The finite-class simultaneous guarantee in
\cref{thm:p15-conc-finite-learning} accounts for that search under its stated
sampling assumptions.

\subsection{Association, intervention and nonidentification}
\label{sec:p14-counter-causal}

Take mutually independent disturbance bits \(U_Z,U_X,U_Y\), with \(U_Z\)
fair and each other bit equal to one with probability \(1/4\). Specify the
acyclic structural equations \(Z=U_Z\), \(X=Z\oplus U_X\),
\(Y=Z\oplus U_Y\), where \(\oplus\)
is exclusive-or. The complete finite model and its normalization are
\cref{def:p16-confounded-binary-model}. Marginalizing \(Z\) gives
\[
\begin{array}{c|cc}
 &Y=0&Y=1\\\hline
 X=0&5/16&3/16\\
 X=1&3/16&5/16
\end{array}
\]
Hence \(P(Y=1\mid X=1)=5/8\),
\(P(Y=1\mid X=0)=3/8\). Replacing the equation for \(X\) by
\(X=x\) leaves the equation and law for \(Y\) unchanged, so
\(P(Y=1\mid\operatorname{do}(X=x))=1/2\) for either \(x\).
The association difference is \(1/4\) while the intervention effect is zero.
These are the results proved in \cref{thm:p16-causal-confounded-example}
under the equation-replacement definition of
\cref{def:p16-causal-intervention}\cite{Pearl2009CausalOverview}.

For nonidentification, retain the two-variable allowed graph \(X\to Y\).
The first model may express the preceding common cause as dependent
disturbances; its \(Y\) equation ignores the allowed \(X\) input.
The second model uses a fair \(X\), an independent bit \(D\) with
\(P(D=1)=3/8\), and equation \(Y=X\oplus D\). It gives exactly the
same four observational masses, but after \(\operatorname{do}(X=1)\)
its \(Y=1\) probability is \(5/8\), rather than \(1/2\).
The full disturbance spaces and functions are in
\cref{thm:p16-causal-nonidentification,annex:p16.causal-nonidentification}.
The compared class allows dependent disturbances; the second model's
independence does not impose that condition on the whole class.

The same theorem gives a separate positive-support failure even in the
independent-disturbance class. On a singleton exogenous space take
\(X=0,Y=0\) in one model and \(X=0,Y=X\) in another, keeping the allowed
graph \(X\to Y\). Both observations are concentrated at \((0,0)\).
After setting \(X=1\), the first has \(Y=0\), the second \(Y=1\).
Observational conditioning at \(X=1\) is undefined; filling its zero
denominator arbitrarily would not identify either structural response.

\subsection{A zero derivative with nonzero output uncertainty}
\label{sec:p14-counter-linearization}

For \(X\) uniform on \([-a,a]\), \(a>0\), consider output \(Y=X^2\).
The derivative at the input mean zero is zero. Linear propagation there
would give output estimate and variance zero. The exact integrals instead
give
\[
 E(Y)=a^2/3,\qquad E(Y^2)=a^4/5,\qquad
 \operatorname{Var}(Y)=4a^4/45.
\]
This is \cref{def:p15-measurement-zero-gradient}. The approximation has
discarded the first nonzero term rather than demonstrated absence of
uncertainty. Conversely, a variable taking only \(-a,a\) can be
nondegenerate while its square is constant; nondegeneracy alone does not
supply the output conclusion. Controlled nonlinear propagation requires
the hypotheses and remainder bounds in
\cref{thm:p15-measurement-taylor-bound}.

\subsection{Ordinally identical prizes with different risk choices}
\label{sec:p14-counter-utility}

Let sure prizes satisfy \(\ell\prec c\prec h\). Compare a half chance
of each extreme with the sure middle prize:
\(p=\tfrac12\delta_\ell+\tfrac12\delta_h\), \(q=\delta_c\).
With \(u(\ell)=0,u(h)=1,u(c)=2/5\), expected utility selects
\(p\), since \(1/2>2/5\). With the same certain-prize ordering but
\(u(c)=3/5\), it selects \(q\). Thus the ordinal ordering of sure
consequences does not fix lottery preferences.

The same worked example \cref{def:p17-vnm-worked-choice} applies the
strictly increasing transformation \(t\mapsto\sqrt t\) to the first
consequence utilities. It leaves the extreme lottery's score \(1/2\)
but changes the middle prize's score to \(\sqrt{2/5}>1/2\), reversing
the comparison. Transforming the already averaged scores instead preserves
their ranking. The permissible consequence-scale transformations are positive
affine, under \cref{thm:p17-vnm-affine-uniqueness}; the two operations
must not be conflated \cite{Rubinstein2015LectureNotes}.

For ambiguity, \cref{thm:p17-ellsberg-two-urns} gives a different
representation failure. With common favorable and unfavorable prizes,
preferring a known fair urn for both red and black bets over an ambiguous
urn would require one additive subjective prior to make both ambiguous
color probabilities less than \(1/2\). Their sum must be one, which
excludes that combination. The conclusion concerns the common one-prior
representation; it does not assert the nonexistence of another decision
model \cite{Ellsberg1961Ambiguity}.

\subsection{Coherent individuals and an incoherent majority}
\label{sec:p14-counter-aggregation}

Use the agenda containing \(p,q,p\land q\) and their negations, with the
three complete truth judgments below. A zero in the conjunction column means
acceptance of its negation.
\[
\begin{array}{c|ccc}
 \text{individual}&p&q&p\land q\\\hline
 1&1&1&1\\
 2&1&0&0\\
 3&0&1&0\\\hline
 \text{strict majority}&1&1&0
\end{array}
\]
Every individual row is satisfiable. Counting assent separately gives
\(\{p,q,\neg(p\land q)\}\), an unsatisfiable collective set. The
profile and complete closure calculation are
\cref{def:p17-ja-doctrinal-profile,thm:p17-ja-paradox}
\cite{ListPettit2001JudgmentPreprint}. The majority set is complete and
weakly consistent in the entry's sense of avoiding each proposition together
with its negation. It is neither jointly satisfiable nor deductively closed.
Thus those weaker properties are not substitutes for joint coherence.

For three premises, rows \((1,0,1),(0,1,1),(1,1,0)\) give majority
support to each premise while every individual rejects the conjunction.
This variant in the same theorem shows why unanimous rejection of a
conclusion need not match separate majority judgments of its premises.
The impossibility result \cref{thm:p17-ja-impossibility} has additional
domain and rule axioms; this one profile illustrates the defect of majority
without proving that every aggregation rule must fail.

\section{Links between each example and the hypothesis it tests}
\label{sec:p14-counter-hypothesis-links}

The following lookup keeps the falsified inference or removed assumption
beside the usable witness and governing result. A listed failure does not
retire the result that held under its original hypotheses.

\begin{longtable}{@{}p{.25\textwidth}p{.30\textwidth}p{.35\textwidth}@{}}
\caption{Links between each example and the hypothesis it tests}\\
\toprule
Proposed transfer or inference & Witness or boundary here & Governing destination \\
\midrule
\endfirsthead
\caption[]{Links between each example and the hypothesis it tests (continued)}\\
\toprule
Proposed transfer or inference & Witness or boundary here & Governing destination \\
\midrule
\endhead
\midrule\multicolumn{3}{r}{\emph{continued on next page}}\\
\endfoot
\endlastfoot
Converse or inverse of a material implication &
\cref{sec:p14-counter-boolean} &
\cref{def:p18-rd-consequent,thm:p18-rd-independent-cases}. \\
Universal information supplies a predicate witness &
\cref{sec:p14-counter-empty-predicate} &
\cref{def:p02-syllogistic-empty-term,tab:p02-syllogistic-lookup}. \\
Classical excluded middle or double-negation elimination is intuitionistically valid &
\cref{sec:p14-counter-ipl} &
\cref{def:p08-ipl-forcing,fig:p08-ipl-excluded-middle}. \\
Every inconsistent premise pair explodes in LP &
\cref{sec:p14-counter-lp} & \cref{def:p08-lp-nonexplosion}. \\
Seriality, reflexivity, symmetry and transitivity can substitute for one another &
\cref{sec:p14-counter-modal} &
\cref{thm:p08-mce-separations,annex:p8.mce.10}. \\
Universal introduction may retain its eigenvariable in an antecedent &
\cref{sec:p14-counter-eigenvariable} &
\cref{def:p07-classical-eigenvariable-case,sheet:lk-defneg}. \\
A local premise may be necessitated &
\cref{sec:p14-counter-local-necessitation} &
\cref{def:p08-mce-language,thm:p08-mce-separations}. \\
An arbitrary defense fixed point is complete; stable extensions always exist &
\cref{sec:p14-counter-extensions} &
\cref{thm:p18-example-extensions,def:p18-acceptance-queries}. \\
First-order compactness transfers to finite models &
\cref{sec:p14-counter-finite-compactness} &
\cref{def:p03-finite-compactness-failure,thm:fol-compactness}. \\
Full second-order compactness follows from first-order compactness &
\cref{sec:p14-counter-second-order-compactness} &
\cref{thm:p02-pa2-categoricity,thm:p02-sol-full-noncompact}. \\
One bounded-round comparison proves unbounded inexpressibility &
\cref{sec:p14-counter-definability} &
\cref{thm:p03-ef-inexpressibility,def:p03-ef-parity}. \\
Recognition of a program-domain set supplies a total decider &
\cref{sec:p14-counter-decidability} &
Section~\ref{sec:universal-reductions-01} and
\cref{thm:halting-undecidable}. \\
Pairwise independence is mutual independence, or survives selection &
\cref{sec:p14-counter-independence} &
\cref{def:p15-pairwise-not-mutual,thm:p15-conc-dependence-boundaries}. \\
An alert likelihood is its hypothesis posterior &
\cref{sec:p14-counter-base-rate} &
\cref{def:p15-alert-bayes-case,thm:p15-finite-bayes}. \\
Posterior interval content supplies uniform fixed-parameter coverage &
\cref{sec:p14-counter-coverage} & \cref{thm:p15-bayes-coverage}. \\
Selected zero training error supplies a small population error &
\cref{sec:p14-counter-selection} &
\cref{thm:p15-conc-selection-example,thm:p15-conc-finite-learning}. \\
Conditioning on treatment equals intervention; an observational law always identifies an effect &
\cref{sec:p14-counter-causal} &
\cref{thm:p16-causal-confounded-example,thm:p16-causal-nonidentification}. \\
A zero linear sensitivity implies zero output uncertainty &
\cref{sec:p14-counter-linearization} &
\cref{def:p15-measurement-zero-gradient,thm:p15-measurement-taylor-bound}. \\
Ordinal consequence utilities determine risky choice or tolerate every increasing rescaling &
\cref{sec:p14-counter-utility} &
\cref{def:p17-vnm-worked-choice,thm:p17-vnm-affine-uniqueness}. \\
Separate majorities preserve individuals' joint coherence &
\cref{sec:p14-counter-aggregation} &
\cref{def:p17-ja-doctrinal-profile,thm:p17-ja-paradox}. \\
\bottomrule
\end{longtable}

To reuse a witness, first match the original assertion's quantifiers, domain
and semantics. Then check that the proposed object satisfies its premises
and every retained hypothesis. Identify the specific false conclusion or
unequal target values. A witness violating a retained assumption is outside
the theorem's scope, even if it is mathematically interesting. For a
nonexistence claim over procedures, identify the reduction or diagonal
construction that works against every eligible candidate.

The theorem links preserve their stated proof forms. The modal, program-domain,
statistical, causal, utility and judgment-aggregation results above link to
their existing proofs. The EF theorem, its general inexpressibility criterion
and the graph-connectivity strategy retain cited proof detail; the displayed
full-second-order effective-completeness reduction retains its labelled sketch.
The elementary model evaluations and worked examples retain their own
definitions and calculations. A collected set of examples supports those
particular failures; it does not establish a census of every failure or an
automatic countermodel construction for every logic.
